
This section addresses the problem of designing a convergent position
and attitude controller for the vehicle.

The dynamical model of the vehicle, in the absence of gravity, can be
derived from the Newton and Euler equations of motion. Let us denote the position and velocity of the
body frame $\mathcal{B}$ with respect to the inertial frame
$\mathcal{I}$ as $\bar{x}$ and $\bar{v}$, the rotation matrix of frame
$\mathcal{B}$ with respect to $\mathcal{I}$ as $\mathbf{R}$, and the
angular velocity of the vehicle in the body frame $\mathcal{B}$ as
$\bar{\omega}$. Then,
\begin{equation}
  \label{eq:dynsys}
  \left\{
    \begin{aligned}
      \dot{\bar{x}} &= \bar{v} \\
      m \dot{\bar{v}} &= \mathbf{R} \bar{F} \\
      \dot{\mathbf{R}} &= \mathbf{R}\,\mathbf{S}(\bar{\omega}) \\
      \mathbf{J} \dot{\bar{\omega}} &= \bar{M} - \bar{\omega}\times\mathbf{J} \bar{\omega}
    \end{aligned}
  \right.
\end{equation}
where the constants $m$ and $\mathbf{J}$
are the vehicle's mass and moment of inertia, while
$\mathbf{S}(\bar{\omega})$ is the skew-symmetric matrix defined by
\begin{equation}
  \label{eq:S}
  \mathbf{S}(\bar{\omega}) =
  \left[
    \begin{array}{ccc}
      0 & -\omega_z & \omega_y \\
      \omega_z & 0 & -\omega_x \\
      -\omega_y & \omega_x & 0 \\
    \end{array}
  \right]
\end{equation}
for $\bar{\omega}=[\omega_x\:\omega_y\:\omega_z]^T$.

The approach used for the motion control of the vehicle exploits its
holonomic design by
decoupling the translational and rotational modes. To do so, we first
apply feedback linearisation~\cite{sastry99} to the translational part
of~\eqref{eq:dynsys}:
\begin{equation}
  \label{eq:fctrl}
    \left\{
    \begin{split}
      \dot{\bar{x}} &= \bar{v} \\
      \dot{\bar{v}} &= \bar{p} \\
      \bar{F} &= m \mathbf{R}^T \bar{p}
    \end{split}
  \right.
\end{equation}
and then design a feedback controller for $\bar{p}$. Since this
dynamical system is diagonal and second order, a PD controller is
enough to ensure exponential
convergence:
\begin{equation}
  \label{eq:pctrl}
  \begin{split}
    \bar{e}_x &= \bar{x} - \bar{x}_d \\
    \bar{e}_v &= \bar{v} - \bar{v}_d \\
    \bar{p} &= -k_x \bar{e}_x -k_v \bar{e}_v \\
  \end{split}
\end{equation}
where $\bar{x}_d$ and $\bar{v}_d$ are the desired position and
velocity vectors in the inertial frame~$\mathcal{I}$, and $k_x$ and
$k_v$ are the proportional and derivative gains of the PD controller.

For the attitude control we follow the exponentially convergent $SO(3)$ controller
proposed in~\cite{lee12}:
\begin{equation}
  \label{eq:mctrl}
  \begin{aligned}
    \bar{e}_R &= \frac{1}{2 \sqrt{1+tr[\mathbf{R}^T_d\mathbf{R}]} } S^{-1}( \mathbf{R}^T_d \mathbf{R} -
    \mathbf{R}^T \mathbf{R}_d ) \\
    \bar{e}_\omega &= \bar{\omega} - \mathbf{R}^T
    \mathbf{R}_d\,\bar{\omega}_d \\
    \bar{M} &= -k_R\,\bar{e}_R - k_\omega\,\bar{e}_\omega \\
    &+\mathbf{S}(\mathbf{R}^T\mathbf{R}_d\bar{\omega}_d)\mathbf{J}\mathbf{R}^T\mathbf{R}_d\bar{\omega}_d 
    +\mathbf{J}\mathbf{R}^T\mathbf{R}_d\dot{\bar{\omega}}_d \\
  \end{aligned}
\end{equation}
where $\mathbf{R}_d$ is the rotation matrix of the desired attitude,
$\bar{\omega}_d$ is the desired angular velocity vector, with $k_R$ and
$k_\omega$ as controller gains. The $S^{-1}$ function performs the inverse
operation as the one defined in~\eqref{eq:S}, that is, recovers
the vector from a given skew-symmetric matrix.


%%% Local Variables:
%%% mode: latex
%%% TeX-master: "main"
%%% End:
